% Build: lualatex mrs-format.tex (twice)
\documentclass[11pt, a4paper]{article}

\usepackage[margin=30mm]{geometry}
\usepackage{amsmath}
\usepackage{fontspec}
\usepackage{microtype}
\usepackage{parskip}
\usepackage{booktabs}
\usepackage{tabularx}
\usepackage{caption}
\usepackage{tikz}
\usetikzlibrary{decorations.pathreplacing}
\usepackage[hidelinks]{hyperref}

\newcommand{\code}[1]{\texttt{#1}}
\newcommand{\debrunner}[1]{(Debrunner 2011, p.~#1)}
\newcolumntype{L}{>{\raggedright\arraybackslash}X}
\setlength{\emergencystretch}{2em}

% Tables and figures stay where they are written, caption attached.
\newenvironment{block}
  {\par\medskip\noindent\begin{minipage}{\textwidth}\centering}
  {\end{minipage}\par\medskip}

\title{MRS and MRSC File Format Specification}
\author{Niels Pfeffer}
\date{15 September 2026}

\begin{document}

\maketitle

\section{Overview}

The roll scanner of Bern University of Applied Sciences (BFH) records a roll with two line cameras, a monochrome camera in transmitted light and an RGB camera in reflected light.
An encoder wheel on the paper triggers both cameras every 0.2\,mm \debrunner{39}.
The data of the two cameras are written to \code{*.mrs} (transmitted light) and \code{*.mrsc} (reflected light) \debrunner{51}.
A reference implementation is \code{mrs2image.py} in \url{https://github.com/pfefferniels/rollscan2image}.

Offsets and lengths are in bytes. Numbers in header fields are decimal ASCII, zero-padded to ten digits.

\section{Container}

Both formats consist of a 124-byte header, raster data in chunks, each preceded by a marker, and an end marker.
MRS files end with a metadata trailer.

\begin{block}
\newcommand{\segment}[4][white]{%
  \filldraw[fill=#1] (#2,0) rectangle ++(#3,9);
  \node[align=center] at ({#2+#3/2},4.5) {#4};}
\begin{tikzpicture}[x=1mm, y=1mm, font=\scriptsize]
  \segment{0}{20}{header\\124}
  \segment[black!12]{20}{8}{$M_0$}
  \segment{28}{18}{chunk 0}
  \segment[black!12]{46}{8}{$M_1$}
  \segment{54}{18}{chunk 1}
  \node at (75,4.5) {\dots};
  \segment[black!12]{78}{10}{$M_{b-1}$}
  \segment{88}{18}{chunk $b-1$}
  \segment[black!12]{106}{7}{$E$}
  \segment{113}{12}{Def.\\size}
  \segment{125}{16}{INI}
  \draw[decorate, decoration={brace, mirror, amplitude=2mm}] (20,-1.5) -- (113,-1.5)
    node[midway, below=2.5mm] {$d$};
  \draw[decorate, decoration={brace, mirror, amplitude=2mm}] (113,-1.5) -- (141,-1.5)
    node[midway, below=2.5mm] {MRS only};
\end{tikzpicture}
\captionof{figure}{Container layout, not to scale. $M_i$ is the marker of chunk $i$, $E$ the end marker.}
\label{fig:container}
\end{block}

\begin{block}
\begin{tabularx}{\textwidth}{@{}llL@{}}
\toprule
Offset & Length & Content \\
\midrule
0 & 30 & \code{Headsize in Bytes = }$h$ \\
30 & 64 & \code{␣Date = }\textit{DD.MM.YYYY}\code{, MRS }\textit{kk}\code{ Nr = V.}\textit{x.y}\code{, Scan Number = }$s$\code{;} \\
94 & 30 & \code{Datasize in Bytes = }$d$ \\
124 & 10 & marker \code{\$000000000} \\
134 & $kL$ & chunk 0 \\
 & 10 & marker \code{\$000000001} \\
 & $kL$ & chunk 1, and so on \\
$114 + d$ & 10 & end marker \code{\$999999999} \\
$124 + d$ & 30 & MRS only: \code{Def.size in Bytes = }$m$ \\
$154 + d$ & $m$ & MRS only: metadata \\
\bottomrule
\end{tabularx}
\captionof{table}{Byte layout. ␣ is a space.}
\label{tab:container}
\end{block}

\begin{itemize}
\item The header fields follow each other without separators.
\item $h$ is the length of the middle field, always 64.
\item \textit{kk} is \code{SW} in MRS and \code{Co} in MRSC, \textit{x.y} is the format version and $s$ the scan number.
\item $d$ is the number of bytes from offset 124 to the end of the end marker.
\item A marker is \code{\$} followed by the chunk number in nine digits, starting at 0.
\item A chunk holds $k$ lines of $L$ bytes each (Table~\ref{tab:rasters}).
\end{itemize}

The number of chunks $b$ and the number of lines $z$ follow from $d$:
\[
b = \left\lceil \frac{d - 10}{kL + 10} \right\rceil
\qquad
z = \frac{d - 10\,(b + 1)}{L}
\]

\section{Rasters}

\begin{block}
\begin{tabular}{@{}lll@{}}
\toprule
 & MRS & MRSC \\
\midrule
Camera & transmitted light, monochrome & reflected light, RGB \\
\textit{kk} & \code{SW} & \code{Co} \\
Pixels per line $n$ & 2048 & 2098 \\
Bytes per pixel & 1 (grey) & 3 (red, green, blue) \\
Bytes per line $L$ & 2048 & 6294 \\
Lines per chunk $k$ & 1000 & 1 \\
Line spacing & 0.2\,mm & 0.2\,mm \\
Pixel spacing, nominal & 0.22\,mm & 0.21\,mm \\
Holes & bright & colour of the backing \\
Metadata trailer & yes & no \\
\bottomrule
\end{tabular}
\captionof{table}{The two rasters. Nominal pixel spacing after \debrunner{40}.}
\label{tab:rasters}
\end{block}

Values are unsigned 8-bit integers.
Line 0 is the start of the scan. A line spans the camera's field of view, which extends beyond both paper edges.
Both cameras scan across the roll in the same direction. The two rasters are not registered to each other.
The pixel spacing across the roll is not stored. It is determined per scan from \code{Roll Width} or from the track grid of the roll.

\section{Metadata trailer}

After the end marker, MRS files contain the 30-byte field \code{Def.size in Bytes = }$m$, followed by $m$ bytes of INI text up to the end of the file.
Lines have the form \textit{key}\code{=}\textit{value} and are grouped under section names in square brackets. Line endings are CR\,LF, the encoding is Latin-1. Unset values are empty.

The keys are listed in Table~\ref{tab:trailer}. $N$ is the value of \code{Number Of Tracks}.
\code{Roll Width}, \code{First Track} and \code{Last Track} are in mm, the track positions measured from the same paper edge.
\code{Play Speed} is in mm/s.

\begin{block}
\begin{tabularx}{\textwidth}{@{}lL@{}}
\toprule
Section & Keys \\
\midrule
\code{[Soundtrack 1]} & \code{Soundtrack Name}, \code{Composer First Name}, \code{Composer Last Name}, \code{Arranger First Name}, \code{Arranger Last Name}, \code{Interpret First Name}, \code{Interpret Last Name} \\
\addlinespace
\code{[Rolltype]} & \code{Name}, \code{Track Shifts}, \code{Control Track Position}, \code{Number Of Tracks}, \code{Roll Width}, \code{First Track}, \code{Last Track}, \code{Track Alignement}, \code{Track Width}, \code{Trackgap}, \code{Min Track Length}, \code{Track 1} \dots\ \code{Track }$N$ \\
\addlinespace
\code{[Definition]} & \code{Title}, \code{Recording Year}, \code{Recording Location}, \code{Producer}, \code{Running Time}, \code{Play Speed}, \code{Notes}, \code{Catalogue Number}, \code{Catalogue Number 2}, \code{Is Replica}, \code{Paper Type}, \code{Roll Length}, \code{Box State} \\
\bottomrule
\end{tabularx}
\captionof{table}{Sections and keys of the metadata trailer, in file order.}
\label{tab:trailer}
\end{block}

\section{Settings file}

The recording software writes a CSV file with the base name of the MRSC file.
Every field is terminated by a semicolon. The first line is \code{DateTime;Group;Variable;Value;Unit;}.
Group \code{Images.settings} holds settings without unit, group \code{Images} holds measurements with the unit in square brackets.

\begin{block}
\begin{tabularx}{\textwidth}{@{}lL@{}}
\toprule
Variable & Content \\
\midrule
\code{MRSC\_Brightness\_Correction\_Enable} & brightness correction (\code{True}, \code{False}) \\
\code{MRSC\_Brightness\_Courve\_A}, \code{\_B}, \code{\_C} & coefficients $A$, $B$, $C$ of the brightness curve \\
\code{MRSC\_SpatialCorrOfOneColor} & colour row offset $\delta$ in lines \\
\code{MRSC\_DetectPencil}$j$\code{\_Enable} & pencil detection in range $j$ (\code{True}, \code{False}) \\
\code{MRSC\_DetectPencil}$j$\code{\_minValues} & lower colour of range $j$ \\
\code{MRSC\_DetectPencil}$j$\code{\_maxValues} & upper colour of range $j$ \\
\code{MRSC\_DetectTape}$j$\code{\_Enable} & tape detection in range $j$ (\code{True}, \code{False}) \\
\code{MRSC\_DetectTape}$j$\code{\_minValues} & lower colour of range $j$ \\
\code{MRSC\_DetectTape}$j$\code{\_maxValues} & upper colour of range $j$ \\
\code{MRSC\_DetectedPixels\_Overwrite} & overwrite detected pixels (\code{True}, \code{False}) \\
\code{MRSC\_Length} & length of the scan in m \\
\code{MRSC\_PensilPixels} & pencil pixels in permille \\
\code{MRSC\_TapePixels} & tape pixels in permille \\
\bottomrule
\end{tabularx}
\captionof{table}{Variables of the settings file, $j \in \{1, 2\}$. Colours are written as \code{Color [A=}$a$\code{, R=}$r$\code{, G=}$g$\code{, B=}$b$\code{]}.}
\label{tab:settings}
\end{block}

\section{Colour image}

The PNG delivered by the scanner is computed from the MRSC raster by colour registration, brightness correction and mirroring.
Scan lines become columns, pixels become rows:
\[
p(x,\, n - 1 - s,\, i) = \min\left(255,\ \left\lfloor r(x + \delta_i,\, s,\, i) \cdot \frac{c_{\max}}{c(s)} \right\rfloor\right)
\]
\begin{itemize}
\item $r(x, s, i)$ is the MRSC value in line $x$, pixel $s$ and channel $i \in \{\mathrm{R}, \mathrm{G}, \mathrm{B}\}$.
\item $\delta_{\mathrm{R}} = 2\delta$, $\delta_{\mathrm{G}} = \delta$, $\delta_{\mathrm{B}} = 0$, with $\delta$ = \code{MRSC\_SpatialCorrOfOneColor}.
\item $c(s) = As^2 + Bs + C$, and $c_{\max}$ is the maximum of $c(s)$ for $0 \le s < n$.
\item If \code{MRSC\_Brightness\_Correction\_Enable} is \code{False}, $c_{\max}/c(s) = 1$.
\item The image has $z - 2\delta$ columns and $n$ rows.
\end{itemize}

\section*{References}

Debrunner, Daniel. 2011. “Die Entwicklung des Musikrollenscanners der Berner Fachhochschule – aus Musikrollenbildern wird Musik – die elektronische Steuerung der Welte-Philharmonie-Orgel.” In \textit{Wie von Geisterhand. Aus Seewen in die Welt. 100 Jahre Welte-Philharmonie-Orgel}, 35–63. Seewen.

\end{document}
